考虑自旋极化的一维开边界Kitaev超导链，费米算符 \(c_j\) 满足反对易关系，哈密顿量
\[
H=-\mu\sum_{j=1}^L(c_j^\dagger c_j-\tfrac12)
-t\sum_{j=1}^{L-1}(c_j^\dagger c_{j+1}+\mathrm{h.c.})
+\Delta\sum_{j=1}^{L-1}(c_jc_{j+1}+\mathrm{h.c.}),
\]
\(t>0,\Delta>0,\mu\in\mathbb R\)。定义 Majorana \(\gamma_{Aj}=c_j+c_j^\dagger\)、\(\gamma_{Bj}=-i(c_j-c_j^\dagger)\)。
(1) 对相应周期链写出BdG动量空间哈密顿量及两条能带，确定 \(\Delta\ne0\) 的体能隙关闭条件。
(2) 对开链特殊点 \(\mu=0,\Delta=t\)，用Majorana算符精确重写哈密顿量，找出两端与 \(H\) 对易的局域零模及其组合的非局域费米占据。
(3) 用 \(q(k)=-\mu-2t\cos k+2i\Delta\sin k\) 的复平面轨迹判定拓扑非平庸区间，说明该绕数为何在开链表现为边界零模。
(4) 解释哈密顿量保留总费米奇偶而不保留总粒子数；研究有限长度两端零模耦合时能隙分裂及 \(t,\Delta,\mu\) 的边界条件，避免把有限链的简并当作在任意参数下严格精确。
